\section{Nghiên cứu tổng quan và cơ sở lý thuyết}

\subsection{Quy trình phát triển phần mềm và công cụ quản lý dự án}
Do đặc thù của đề tài kết hợp đồng thời phát triển trò chơi 3D, thử nghiệm mạng trực tuyến và nghiên cứu nhiều kỹ thuật học máy có độ bất định cao, nhóm áp dụng quy trình phát triển lặp tăng trưởng theo tinh thần \textbf{Agile} với chu kỳ kiểm nghiệm theo tuần (15 tuần cho Giai đoạn 1). Khác với mô hình thác nước đòi hỏi khóa cứng toàn bộ đặc tả từ đầu, cách tiếp cận lặp cho phép nhóm nhanh chóng điều chỉnh tham số cân bằng điểm hành động và hợp thể, chuyển đổi giải pháp mạng từ Unity NGO sang Photon Fusion dựa trên số liệu đo đạc thực tế, và cải tiến liên tục hai hướng tác nhân AI là Hybrid Bot và LLM Bot.

Các công cụ hỗ trợ quản lý và triển khai gồm:
\begin{itemize}
    \item \textbf{Git và GitHub:} Quản lý phiên bản mã nguồn phân nhánh độc lập theo vai trò giữa hai thành viên, giúp theo dõi lịch sử thay đổi xuyên suốt giai đoạn 1 và hạn chế xung đột trên các tài nguyên đồ họa Unity.
    \item \textbf{Unity 6:} Được lựa chọn so với Unreal Engine và Godot nhờ cân bằng tốt giữa năng suất lập trình C\# (thuận lợi cho việc tách riêng lõi luật chơi không phụ thuộc đồ họa), hệ sinh thái gói chính chủ \textbf{Unity ML-Agents} \cite{mlagents}, khả năng tích hợp mượt mà với \textbf{Photon Fusion 2} \cite{photon_fusion} và trình kết xuất đồ họa nhẹ phù hợp với phần cứng máy tính tầm trung (bộ nhớ đồ họa 6 GB) \cite{unity6}.
    \item \textbf{Python, PyTorch, ONNX Runtime và Blender 4.2 (MPFB v2):} Hệ sinh thái tiêu chuẩn để huấn luyện các mô hình học máy cho Hybrid Bot, chạy vòng lặp tiến hóa luật cho LLM Bot bằng Gemini API và xây dựng quy trình tái tạo khuôn mặt cùng tóc 3D cho Smart Customization \cite{mpfb2}.
\end{itemize}

\subsection{Các nguyên tắc và mẫu thiết kế kiến trúc áp dụng}
Để giữ cấu trúc phần mềm gọn gàng và dễ mở rộng trong suốt hai giai đoạn thực hiện, hệ thống vận dụng có chọn lọc năm mẫu thiết kế kiến trúc sau:
\begin{itemize}
    \item \textbf{Kiến trúc lục giác tách biệt lõi và giao diện \cite{cockburn}:} Đặt toàn bộ thực thể và quy tắc cờ trong vùng lõi độc lập khỏi đồ họa 3D. Mọi tương tác vào ra của ván đấu đều đi qua hai cổng giao tiếp chuẩn: một cổng tiếp nhận nước đi (từ người chơi, từ mạng hoặc từ AI) và một cổng phát sự kiện diễn biến ván đấu cho phần hiển thị hoặc bộ ghi nhận thực nghiệm.
    \item \textbf{Cô lập dịch vụ theo từng bàn đấu:} Thay vì dùng biến dùng chung toàn cục dễ gây xung đột khi chạy song song nhiều bàn cờ huấn luyện AI, mỗi bàn đấu sở hữu một vùng quản lý dịch vụ riêng được cấu hình độc lập tùy theo chế độ chơi (đấu với máy, đấu hai người hay chạy mô phỏng huấn luyện).
    \item \textbf{Chuẩn hóa lệnh hành động và gom nhóm sự kiện:} Mọi nước đi đều được đóng gói thành một cấu trúc lệnh thống nhất; sau khi lõi luật xử lý xong sẽ phát ra một gói sự kiện theo đúng trình tự xảy ra, giúp tầng hiển thị diễn hoạt tuần tự còn các bộ theo dõi huấn luyện vẫn lắng nghe độc lập.
    \item \textbf{Máy trạng thái hữu hạn cho hoạt họa quân cờ:} Quản lý rõ ràng 5 trạng thái hoạt họa của từng quân cờ 3D (đứng chờ, di chuyển, tấn công, bị đánh trúng và bị tiêu diệt), giúp mã xử lý hoạt họa gọn gàng và dễ mở rộng thêm mẫu quân cờ mới.
    \item \textbf{Mẫu thiết kế khởi tạo theo loại bản tin và bọc bộ điều khiển trong tầng mạng:} Áp dụng cho phân hệ Multiplayer để tự động phân tuyến gói tin nhị phân theo mã định danh và bổ sung tính năng gửi nước đi lên mạng mà không cần sửa đổi bộ điều khiển chuột cục bộ.
\end{itemize}

\subsection{Cơ sở lý thuyết và đánh giá giải pháp Networking}
\subsubsection{So sánh các giải pháp Networking cho Unity}
Trong giai đoạn phác thảo đề cương ban đầu, nhóm từng cân nhắc Unity Netcode for GameObjects (NGO). Tuy nhiên, khi đi vào phân tích kiến trúc cụ thể của \textit{Convergence Chess} (lõi luật tách rời \texttt{GameObject}, lối chơi theo lượt cần độ tin cậy tuyệt đối và không muốn phụ thuộc vào hạ tầng cấu hình Relay/Lobby phức tạp), nhóm đã khảo sát và so sánh ba giải pháp phổ biến trong hệ sinh thái Unity như Bảng \ref{tab:net_solutions}.

\begin{table}[H]
    \centering
    \caption{So sánh các giải pháp Networking cho game chiến thuật theo lượt trên Unity}
    \label{tab:net_solutions}
    \begin{tabularx}{\textwidth}{L{2.5cm} Y Y Y}
        \toprule
        \textbf{Tiêu chí} & \textbf{Unity NGO} & \textbf{Mirror Networking} & \textbf{Photon Fusion 2} \\
        \midrule
        Hạ tầng kết nối & Cần tích hợp thêm Unity Lobby và Unity Relay để vượt NAT & Tự quản lý máy chủ hoặc dùng Transport trung gian (KCP/Telepathy) & Tích hợp sẵn Photon Cloud toàn cầu, chọn Region trực tiếp, không cần dựng server riêng \\
        Mức độ phụ thuộc GameObject & Mạnh (yêu cầu gắn \texttt{NetworkObject} và \texttt{NetworkBehaviour}) & Mạnh (yêu cầu \texttt{NetworkIdentity} trên các đối tượng đồng bộ) & Hỗ trợ cả \texttt{NetworkObject} lẫn truyền luồng byte thuần qua \texttt{SendReliableDataToPlayer} \\
        Phù hợp với Core tách đồ họa & Khó tách rời nếu dùng đồng bộ đối tượng & Cần tự xây lớp cầu nối & Rất cao nhờ cơ chế \texttt{ReliableData} gửi trực tiếp struct/byte payload \\
        Chi phí giai đoạn phát triển & Miễn phí theo hạn mức Unity Gaming Services & Mã nguồn mở miễn phí (tự tốn chi phí thuê VPS/Server) & Miễn phí gói phát triển lên đến 20 CCU trên Photon Cloud \\
        \bottomrule
    \end{tabularx}
\end{table}

\textbf{Kết luận lựa chọn giai đoạn 1:} Nhóm quyết định chuyển từ Unity NGO sang \textbf{Photon Fusion 2} \cite{photon_fusion} (chạy thống nhất ở mô hình \texttt{GameMode.Shared} trên Photon Cloud, trong đó máy chủ phòng là Shared Mode Master Client) và đặc biệt chọn cơ chế truyền luồng dữ liệu tin cậy \textbf{Fusion Reliable Data Stream} thay cho đồng bộ trạng thái quân cờ bằng \texttt{Fusion Network Object}. Lựa chọn này giải quyết đồng thời hai yêu cầu sống còn: (i) cho phép triển khai nhanh chế độ PvP Online đa khu vực mà không phải vận hành máy chủ riêng ở giai đoạn DACN; và (ii) bảo toàn kiến trúc Decoupled Core vì các quân cờ trên bàn không bị biến thành \texttt{NetworkObject}. Số liệu đo đạc thực nghiệm so sánh trực tiếp giữa \texttt{Fusion Network Object} và \texttt{Fusion Reliable Data} sẽ được trình bày chi tiết tại Chương 5 và Chương 6.

\subsection{Cơ sở lý thuyết về trí tuệ nhân tạo trong game chiến thuật}
Bàn cờ $6 \times 6$ của \textit{Convergence Chess} có số ô ít hơn cờ vua $8 \times 8$, nhưng độ phức tạp rẽ nhánh lại mang nét đặc thù do tổ hợp 9 kỹ năng nội tại, kinh tế điểm hành động biến thiên theo khoảng cách di chuyển, cơ chế hợp thể và chuỗi đòn bồi Surge. Vì vậy, để xây dựng hai dòng tác nhân chính là \textbf{Hybrid Bot} và \textbf{LLM Bot}, nhóm nghiên cứu bốn nền tảng lý thuyết bổ trợ cho nhau:

\subsubsection{Tìm kiếm cây trò chơi Minimax với cắt tỉa Alpha-Beta và bảng băm Zobrist}
Thuật toán Minimax tìm nước đi tối ưu trong trò chơi đối kháng hai người theo nguyên lý cực đại hóa lợi thế trong kịch bản đối thủ phản ứng tốt nhất. Để giảm mạnh số trạng thái phải duyệt, nền tảng tìm kiếm của Hybrid Bot kết hợp các kỹ thuật tối ưu kinh điển:
\begin{itemize}
    \item \textbf{Cắt tỉa Alpha-Beta kết hợp sắp xếp nước đi ưu tiên:} Loại bỏ sớm các nhánh chắc chắn không cải thiện kết quả, đồng thời ghi nhớ các nước đi thường gây cắt tỉa tốt và bảng điểm lịch sử cho các nước đi tĩnh để duyệt những nhánh hứa hẹn trước.
    \item \textbf{Bảng hoán vị Zobrist và gia hạn độ sâu khi có đòn bồi Surge:} Mã hóa toàn bộ trạng thái bàn cờ (vị trí quân, vật cản, kỹ năng, trạng thái hợp thể, AP hai bên, số lượt bỏ lượt, trạng thái Surge và lượt đi) thành một mã băm 64-bit để tránh tính lại các thế cờ lặp. Đồng thời, hệ thống gia hạn thêm tối đa $2$ tầng tìm kiếm trên các nhánh kích hoạt đòn bồi Surge nhằm tránh hiện tượng bỏ sót đòn ăn quân liên hoàn ở rìa độ sâu tìm kiếm.
    \item \textbf{Giảm độ sâu thăm dò có kiểm chứng lại theo mạng chính sách:} Khi kết hợp với mạng nơ-ron gợi ý nước đi, bộ tìm kiếm duyệt đủ độ sâu cho nước đi tốt nhất từ bảng hoán vị, nhóm nước đi có xác suất cao và toàn bộ nước ăn quân hoặc hợp thể; với các nước đi tĩnh ít tiềm năng hơn, hệ thống duyệt nông hơn để tiết kiệm thời gian và tự động kiểm chứng lại ở độ sâu đầy đủ nếu phát hiện điểm số vượt trội.
\end{itemize}

\subsubsection{Học bắt chước và Học tăng cường PPO tự đấu (thành phần huấn luyện Hybrid Bot)}
Huấn luyện học tăng cường từ con số không trên trò chơi có tín hiệu thưởng thưa (chỉ phân định thắng thua khi tiêu diệt Vua) thường hội tụ rất chậm. Do đó, mạng chính sách của Hybrid Bot được huấn luyện qua hai bước nối tiếp trên nền tảng \textbf{Unity ML-Agents} \cite{mlagents}:
\begin{itemize}
    \item \textbf{Học bắt chước (Behavioral Cloning và DAgger \cite{dagger}):} Thu thập hàng chục nghìn cặp thế cờ và nước đi mẫu từ bộ tìm kiếm Minimax độ sâu 4 (kết hợp lấy mẫu có độ ngẫu nhiên để đa dạng hóa thế cờ và lấy nhãn tối ưu) nhằm khởi tạo trọng số mạng nơ-ron ban đầu, tránh giai đoạn mò mẫm kém hiệu quả.
    \item \textbf{Học tăng cường tối ưu chính sách cắt ngưỡng (PPO) \cite{ppo}:} Tinh chỉnh chính sách chọn nước đi $\pi_\theta(a|s)$ thông qua quá trình tự đấu bằng hàm mục tiêu cắt ngưỡng, ngăn bước cập nhật quá lớn làm sụp đổ chiến thuật đã học:
    \begin{equation}
        L^{\text{CLIP}}(\theta) = \hat{\mathbb{E}}_t \left[ \min\left( r_t(\theta)\hat{A}_t,\; \text{clip}\left(r_t(\theta), 1-\epsilon, 1+\epsilon\right)\hat{A}_t \right) \right]
    \end{equation}
    trong đó $r_t(\theta) = \frac{\pi_\theta(a_t|s_t)}{\pi_{\theta_{\text{old}}}(a_t|s_t)}$ là tỷ số xác suất giữa chính sách mới và cũ, còn $\hat{A}_t$ là ước lượng mức độ lợi thế của hành động.
    \item \textbf{Định hình phần thưởng theo hàm thế năng \cite{potential_shaping}:} Để cung cấp tín hiệu dẫn hướng đều đặn hơn trong trận mà không làm lệch mục tiêu thắng thua gốc, phần thưởng trung gian được tính từ chênh lệch điểm đánh giá thế cờ của mạng NNUE:
    \begin{equation}
        F(s_t, s_{t+1}) = \lambda \cdot \left( \gamma \Phi_{\text{NNUE}}(s_{t+1}, p) - \Phi_{\text{NNUE}}(s_t, p) \right)
    \end{equation}
\end{itemize}

\subsubsection{Mạng đánh giá nhanh NNUE và biến thể dạng dư (Residual NNUE)}
Lấy cảm hứng từ kiến trúc \textbf{NNUE} trong máy đánh cờ Stockfish \cite{nnue}, thay vì chạy suy luận mạng nơ-ron nặng ở từng nút lá của cây tìm kiếm, nhóm sử dụng một mạng nơ-ron gọn ba tầng ($654 \to 128 \to 64 \to 1$) tận dụng tính thưa của bàn cờ $6 \times 6$: chỉ cộng dồn trọng số tại những ô thực sự có quân cờ hoặc vật cản và bỏ qua toàn bộ ô trống. Ở biến thể \textbf{NNUE dạng dư (Residual NNUE)}, hệ thống giữ lại điểm đánh giá vật chất cơ sở làm mốc neo và chỉ để mạng nơ-ron học phần hiệu chỉnh chiến thuật tinh tế trong biên độ giới hạn $[-3.0, +3.0]$. Cách làm này giúp tránh sai lệch nghiêm trọng về giá trị quân cờ và tăng tốc độ hội tụ rõ rệt.

\subsubsection{Tiến hóa hàm đánh giá biểu tượng bằng Mô hình Ngôn ngữ Lớn (nền tảng LLM Bot)}
Song song với hướng học sâu của Hybrid Bot, nhóm nghiên cứu hướng \textbf{LLM Bot} theo phương pháp \textbf{FunSearch} \cite{funsearch} nhằm tạo ra tri thức chiến thuật có thể đọc hiểu và kiểm chứng được. Hệ thống sử dụng Gemini API để sinh và tiến hóa trực tiếp các đoạn quy tắc đánh giá ngắn gọn, phân loại theo 8 tình huống chiến thuật của ván đấu (khai cuộc chưa hợp thể, trung cuộc đã hợp thể, tàn cuộc ít quân, tranh chấp khiên, Vua bị đe dọa, đua phong cấp Tốt, dự trữ AP cao...). Mỗi quy tắc sinh ra đều trải qua 4 bước sàng lọc tự động về tính hợp lệ, tốc độ, tính đối xứng và tỷ lệ thắng đối kháng trước khi được lưu vào ngân hàng luật của LLM Bot.

\subsection{Cơ sở lý thuyết cho phân hệ tùy biến nhân vật 3D (Smart Customization)}
Bài toán tái tạo nhân vật 3D có sẵn khung xương để hoạt họa từ một bức ảnh chân dung 2D đơn lẻ đòi hỏi giải quyết đồng thời hai bài toán con:
\begin{itemize}
    \item \textbf{Tái tạo hình học khuôn mặt 3D:} Sử dụng \textbf{MediaPipe} \cite{mediapipe} trích xuất 478 điểm mốc giải phẫu trên ảnh 2D, sau đó ánh xạ sang hệ thống thông số biến dạng khuôn mặt của bộ khung chuẩn giải phẫu \textbf{MPFB v2} trên Blender \cite{mpfb2}. Thực nghiệm cho thấy việc dùng trực tiếp chiều sâu ước lượng từ MediaPipe gây sai số lớn (sai số trung bình bình phương đạt $0.1415$ theo khoảng cách hai mắt). Do đó, nhóm áp dụng phương pháp sinh dữ liệu ảnh phối cảnh tổng hợp trực tiếp từ mô hình MPFB v2 bằng \textbf{PyTorch3D} \cite{pytorch3d} để huấn luyện bộ hồi quy gọn trên $25$ thông số hình thái độc lập, kết hợp bước tinh chỉnh tối ưu bình phương tối thiểu có giới hạn.
    \item \textbf{Tái tạo kết cấu bề mặt da, khử bóng và chi tiết nếp nhăn:} Kế thừa ý tưởng từ nghiên cứu \textbf{MeInGame} \cite{meingame}, hệ thống dùng phép chiếu ngược tam giác trên \textbf{PyTorch3D} \cite{pytorch3d} để ánh xạ trực tiếp màu da từ ảnh chân dung sang bản đồ bề mặt của nhân vật 3D. Tiếp đó, hệ thống khử bóng đổ định hướng bằng kỹ thuật đối xứng hai bên khuôn mặt trong không gian màu Lab (tận dụng cấu trúc đối xứng tự nhiên để trung hòa vùng tối lệch hoặc tóc mái che một bên nhưng vẫn giữ được lỗ chân lông, lông mày và râu), kết hợp tách riêng vùng hốc mắt, lọc loại bỏ vệt răng trắng khi ảnh gốc đang cười, hòa trộn biên đa tầng làm mịn vùng chuyển tiếp gò má -- cổ và trích xuất bản đồ chi tiết bề mặt nếp nhăn từ mô hình \textbf{HRN} \cite{hrn}. Hướng khử bóng bằng mô hình khuếch tán \textbf{IC-Light} \cite{iclight} cũng đã được khảo sát thử nghiệm nhưng tạm dừng sử dụng do chi phí tính toán cao và nguy cơ làm biến đổi nét nhận dạng khuôn mặt gốc.
\end{itemize}
